% GENERATED FILE. Edit canonical lessons, then run npm run generate:books.
% canonical-source-hash: fnv1a64:3d519533d87c2dd3
% canonical-lessons: AR-C474-saqr, AR-C474-naama, AR-C474-nawras, AR-C474-risha, AR-C474-tafa

\chapter{Birds}
\label{ch:ar-birds}
\hlchaptermodality{474}

\begin{chapteropening}
\textbf{By the end of this chapter:} \emph{I can say a falcon, an ostrich, a seagull, a feather and to float.}

Complete the last lesson of chapter 474: I can say a falcon, an ostrich, a seagull, a feather and to float.
\end{chapteropening}

\section[\texorpdfstring{\ar{صقر} (ṣaqr)}{صقر (ṣaqr)}]{\ar{صقر} (ṣaqr) — a falcon}

\label{lesson:AR-C474-saqr}



\begin{quote}
Before the new one: say the Arabic for a swamp, then the Arabic for barren.
\end{quote}

\subsection*{You'll want to know: \ar{صقر}}
\textbf{\ar{صقر}} — \emph{ṣaqr} — \textquotedblleft{}a falcon\textquotedblright{}.

\begin{quote}
\ar{طار} \ar{صقر} \ar{فوق} \ar{الوادي}.
\end{quote}

\begin{quote}
\emph{ṭāra ṣaqr fawqa al-wādī.}
\end{quote}

\begin{quote}
A falcon flew over the valley.
\end{quote}

Plural: \ar{صقور}, \emph{ṣuqūr}.

\begin{cousinweb}
\textbf{\ar{صقر}} (\emph{ṣaqr}) is usually traced to the native Arabic root \emph{ṣ-q-r}, \textquotedblleft{}a bird of prey\textquotedblright{}. It is an old native name.
\end{cousinweb}

\begin{grammarlens}[title={What you've built}]
One more word for getting around.
\end{grammarlens}

\subsection*{Your turn}
\begin{itemize}
\raggedright
  \item \emph{Say it:} \emph{ṣaqr}
  \item \emph{Say it:} \emph{ṣaqr}, once more
  \item \emph{Say it:} \emph{ṭāra ṣaqr fawqa al-wādī.} — A falcon flew over the valley.
  \item \emph{Recall:} say the Arabic for a swamp, then the Arabic for barren, then say \emph{ṣaqr} again
\end{itemize}

\begin{tcolorbox}[breakable,colback=teal!4,colframe=teal!35!black,title={Before you move on}]
What does \ar{صقر} mean? (\textquotedblleft{}A falcon\textquotedblright{}.) Say it once more.
\end{tcolorbox}
\section[\texorpdfstring{\ar{نعامة} (naʿāma)}{نعامة (naʿāma)}]{\ar{نعامة} (naʿāma) — an ostrich}

\label{lesson:AR-C474-naama}



\begin{quote}
Before the new one: say the Arabic for barren, then the Arabic for a falcon.
\end{quote}

\subsection*{You'll want to know: \ar{نعامة}}
\textbf{\ar{نعامة}} — \emph{naʿāma} — \textquotedblleft{}an ostrich\textquotedblright{}.

\begin{quote}
\ar{تركض} \ar{النعامة} \ar{بسرعة} \ar{كبيرة}.
\end{quote}

\begin{quote}
\emph{tarkuḍu an-naʿāma bi-surʿa kabīra.}
\end{quote}

\begin{quote}
The ostrich runs very fast.
\end{quote}

\begin{cousinweb}
\textbf{\ar{نعامة}} (\emph{naʿāma}) is an old Arabic name for the ostrich; the origin is not settled.
\end{cousinweb}

\begin{grammarlens}[title={What you've built}]
One more word for getting around.
\end{grammarlens}

\subsection*{Your turn}
\begin{itemize}
\raggedright
  \item \emph{Say it:} \emph{naʿāma}
  \item \emph{Say it:} \emph{naʿāma}, once more
  \item \emph{Say it:} \emph{tarkuḍu an-naʿāma bi-surʿa kabīra.} — The ostrich runs very fast.
  \item \emph{Recall:} say the Arabic for barren, then the Arabic for a falcon, then say \emph{naʿāma} again
\end{itemize}

\begin{tcolorbox}[breakable,colback=teal!4,colframe=teal!35!black,title={Before you move on}]
What does \ar{نعامة} mean? (\textquotedblleft{}An ostrich\textquotedblright{}.) Say it once more.
\end{tcolorbox}
\section[\texorpdfstring{\ar{نورس} (nawras)}{نورس (nawras)}]{\ar{نورس} (nawras) — a seagull}

\label{lesson:AR-C474-nawras}



\begin{quote}
Before the new one: say the Arabic for a falcon, then the Arabic for an ostrich.
\end{quote}

\subsection*{You'll want to know: \ar{نورس}}
\textbf{\ar{نورس}} — \emph{nawras} — \textquotedblleft{}a seagull\textquotedblright{}.

\begin{quote}
\ar{وقف} \ar{نورس} \ar{على} \ar{القارب}.
\end{quote}

\begin{quote}
\emph{waqafa nawras ʿalā al-qārib.}
\end{quote}

\begin{quote}
A seagull stood on the boat.
\end{quote}

\begin{cousinweb}
\textbf{\ar{نورس}} (\emph{nawras}) names the gull; the origin is not settled.
\end{cousinweb}

\begin{grammarlens}[title={What you've built}]
One more word for getting around.
\end{grammarlens}

\subsection*{Your turn}
\begin{itemize}
\raggedright
  \item \emph{Say it:} \emph{nawras}
  \item \emph{Say it:} \emph{nawras}, once more
  \item \emph{Say it:} \emph{waqafa nawras ʿalā al-qārib.} — A seagull stood on the boat.
  \item \emph{Recall:} say the Arabic for a falcon, then the Arabic for an ostrich, then say \emph{nawras} again
\end{itemize}

\begin{tcolorbox}[breakable,colback=teal!4,colframe=teal!35!black,title={Before you move on}]
What does \ar{نورس} mean? (\textquotedblleft{}A seagull\textquotedblright{}.) Say it once more.
\end{tcolorbox}
\section[\texorpdfstring{\ar{ريشة} (rīsha)}{ريشة (rīsha)}]{\ar{ريشة} (rīsha) — a feather; a quill}

\label{lesson:AR-C474-risha}



\begin{quote}
Before the new one: say the Arabic for an ostrich, then the Arabic for a seagull.
\end{quote}

\subsection*{You'll want to know: \ar{ريشة}}
\textbf{\ar{ريشة}} — \emph{rīsha} — \textquotedblleft{}a feather; a quill\textquotedblright{}.

\begin{quote}
\ar{وجد} \ar{الطفل} \ar{ريشة} \ar{بيضاء}.
\end{quote}

\begin{quote}
\emph{wajada aṭ-ṭifl rīsha bayḍāʾ.}
\end{quote}

\begin{quote}
The child found a white feather.
\end{quote}

\begin{cousinweb}
\textbf{\ar{ريشة}} (\emph{rīsha}) is built on the native Arabic root \emph{r-y-sh}, \textquotedblleft{}feathers\textquotedblright{} (Lane). It is one of \emph{rīsh}, \textquotedblleft{}feathers\textquotedblright{}.
\end{cousinweb}

\begin{grammarlens}[title={What you've built}]
One more word for getting around.
\end{grammarlens}

\subsection*{Your turn}
\begin{itemize}
\raggedright
  \item \emph{Say it:} \emph{rīsha}
  \item \emph{Say it:} \emph{rīsha}, once more
  \item \emph{Say it:} \emph{wajada aṭ-ṭifl rīsha bayḍāʾ.} — The child found a white feather.
  \item \emph{Recall:} say the Arabic for an ostrich, then the Arabic for a seagull, then say \emph{rīsha} again
\end{itemize}

\begin{tcolorbox}[breakable,colback=teal!4,colframe=teal!35!black,title={Before you move on}]
What does \ar{ريشة} mean? (\textquotedblleft{}A feather; a quill\textquotedblright{}.) Say it once more.
\end{tcolorbox}
\section[\texorpdfstring{\ar{طفا} (ṭafā)}{طفا (ṭafā)}]{\ar{طفا} (ṭafā) — to float}

\label{lesson:AR-C474-tafa}



\begin{quote}
Before the new one: say the Arabic for a seagull, then the Arabic for a feather.
\end{quote}

\subsection*{You'll want to know: \ar{طفا}}
\textbf{\ar{طفا}} — \emph{ṭafā} — \textquotedblleft{}to float\textquotedblright{}.

\begin{quote}
\ar{طفت} \ar{ريشة} \ar{على} \ar{الماء}.
\end{quote}

\begin{quote}
\emph{ṭafat rīsha ʿalā al-māʾ.}
\end{quote}

\begin{quote}
A feather floated on the water.
\end{quote}

Present: \ar{يطفو}, \emph{yaṭfū}.

\begin{cousinweb}
\textbf{\ar{طفا}} (\emph{ṭafā}) is usually traced to the native Arabic root \emph{ṭ-f-w}, \textquotedblleft{}to rise to the surface of water\textquotedblright{}. It is the plain verb.
\end{cousinweb}

\begin{grammarlens}[title={What you've built}]
One more word for getting around.
\end{grammarlens}

\subsection*{Your turn}
\begin{itemize}
\raggedright
  \item \emph{Say it:} \emph{ṭafā}
  \item \emph{Say it:} \emph{ṭafā}, once more
  \item \emph{Say it:} \emph{ṭafat rīsha ʿalā al-māʾ.} — A feather floated on the water.
  \item \emph{Recall:} say the Arabic for a seagull, then the Arabic for a feather, then say \emph{ṭafā} again
\end{itemize}

\begin{tcolorbox}[breakable,colback=teal!4,colframe=teal!35!black,title={Before you move on}]
What does \ar{طفا} mean? (\textquotedblleft{}To float\textquotedblright{}.) Say it once more.
\end{tcolorbox}
